\documentclass[11pt]{article}
\usepackage[a4paper,margin=1in]{geometry}
\usepackage{fontspec}
\usepackage[boldfont=false]{xeCJK}
\setCJKmainfont{FandolSong-Regular.otf}
\usepackage{amsmath}
\usepackage{amssymb}
\usepackage{array}
\usepackage{booktabs}
\usepackage{graphicx}
\usepackage{adjustbox}
\usepackage{enumitem}
\setlist[itemize]{topsep=2pt,partopsep=0pt,parsep=0pt,itemsep=2pt,leftmargin=1.4em}
\setlist[enumerate]{topsep=2pt,partopsep=0pt,parsep=0pt,itemsep=2pt,leftmargin=1.6em}
\usepackage{parskip}
\usepackage{fvextra}
\fvset{breaklines=true,breakanywhere=true,fontsize=\small}
\setlength{\parindent}{0pt}

\begin{document}

\begin{center}
  {\LARGE\bfseries Ideal gases}\\[0.35em]
  {\large A Level Physics}
\end{center}
\vspace{0.6em}

\section*{The mole}

\textbf{Amount of substance} 物质的量 is an SI base quantity. Its unit is the \textbf{mole} 摩尔 (mol) — one of the seven SI base units, with the kilogram, metre, second, ampere and \textbf{kelvin} 开尔文 from Topic 1.

One mole of any substance has a number of particles equal to the \textbf{Avogadro constant} 阿伏伽德罗常量:

\[
N_{\text{A}} = 6.02 \times 10^{23}\ \text{mol}^{-1}.
\]

A "particle" means whatever you are counting — \textbf{atoms} 原子 for a \textbf{monatomic} 单原子 element like helium, \textbf{molecules} 分子 for $\text{O}_{2}$ or $\text{H}_{2}\text{O}$. Always say what you are counting.

For $n$ moles, the number of particles is $N = n N_{\text{A}}$.

The \textbf{molar mass} 摩尔质量 $M_{\text{m}}$ is the mass of one mole ($\text{kg mol}^{-1}$ or $\text{g mol}^{-1}$). Mass of $n$ moles is $M = n M_{\text{m}}$. Mass of one particle is $m_{0} = M_{\text{m}} / N_{\text{A}}$.

\section*{Equation of state of an ideal gas}

An \textbf{ideal gas} 理想气体 obeys $pV \propto T$ exactly, where $T$ is the \textbf{thermodynamic temperature} 热力学温度.

The \textbf{equation of state} 状态方程 can be written two equal ways:

\[
p V = n R T \qquad\text{or}\qquad p V = N k T.
\]

Here:

\begin{itemize}
\item $p$ — \textbf{pressure} 压强 (Pa).

\item $V$ — \textbf{volume} 体积 (m³).

\item $T$ — thermodynamic temperature in kelvin (never °C).

\item $n$ — number of moles; $N$ — number of molecules.

\item $R$ — \textbf{molar gas constant} 摩尔气体常量, $R = 8.31\ \text{J mol}^{-1}\ \text{K}^{-1}$.

\item $k$ — \textbf{Boltzmann constant} 玻尔兹曼常量, $k = 1.38 \times 10^{-23}\ \text{J K}^{-1}$.

\end{itemize}

Since $N = n N_{\text{A}}$, we get $k = R/N_{\text{A}}$: $k$ is the gas constant \textbf{per molecule}, as $R$ is per mole.

\subsection*{Using the equation of state}

List the variables you have, find the unknown, and choose the form that matches your "amount" (moles → $nRT$; molecules → $NkT$). Always use SI units: Pa, m³, K.

If a fixed amount of gas changes from state 1 to state 2:

\[
\frac{p_{1} V_{1}}{T_{1}} = \frac{p_{2} V_{2}}{T_{2}}.
\]

Special cases:

\begin{itemize}
\item \textbf{constant temperature} (\textbf{Boyle's law} 玻意耳定律): $p_{1} V_{1} = p_{2} V_{2}$.

\item \textbf{constant pressure} (\textbf{Charles's law} 查理定律): $V / T = \text{constant}$.

\item \textbf{constant volume} (\textbf{pressure law} 气体压强定律): $p / T = \text{constant}$.

\end{itemize}

A common mistake is using °C instead of K — $pV \propto T$ only holds with $T$ in kelvin.

\par\medskip
\noindent\begin{minipage}{\linewidth}
\centering
\adjustbox{max width=\linewidth, max totalheight=0.42\textheight}{%
  \includegraphics{/Users/joschka/Desktop/Programming/eXercise/output/eXam/handouts/a_level_physics/assets/15-p271-charles-law-graph.pdf}}\\[0.3em]
\emph{At constant pressure the volume of a gas rises linearly with temperature, reaching zero at absolute zero (Charles's law)}
\end{minipage}\par\medskip

\section*{Kinetic theory of gases}

The \textbf{kinetic theory} 分子动理论 explains a gas's large-scale behaviour from the \textbf{random motion} 无规则运动 of its molecules.

\subsection*{Assumptions}

For an ideal gas:

\begin{enumerate}
\item a \textbf{large number} of identical molecules in continuous random motion.

\item the molecules' own volume is too small to matter compared with the container.

\item the time of each collision is too short to matter compared with the time between collisions.

\item \textbf{intermolecular} 分子间 forces are ignored except during collisions (molecules go in straight lines between them).

\item collisions are \textbf{elastic} — this is an \textbf{elastic collision} 弹性碰撞 (no kinetic energy lost on average).

\item Newton's laws apply.

\end{enumerate}

These assumptions become poor at very high pressure (molecular volume matters) or very low temperature (intermolecular forces matter).

\subsection*{Pressure of a gas — outline of the derivation}

Take a cubic box of side $L$ with $N$ molecules, each of mass $m$. Look at one molecule moving along the $x$-axis with velocity $u_{1}$.

\par\medskip
\noindent\begin{minipage}{\linewidth}
\centering
\adjustbox{max width=\linewidth, max totalheight=0.42\textheight}{%
  \includegraphics{/Users/joschka/Desktop/Programming/eXercise/output/eXam/handouts/a_level_physics/assets/15-p279-molecule-in-box.pdf}}\\[0.3em]
\emph{The pressure derivation considers one molecule's velocity component c\_x normal to a face of a cube of gas}
\end{minipage}\par\medskip

\begin{itemize}
\item \textbf{one collision with the right wall}: velocity reverses to $-u_{1}$, change in \textbf{momentum} 动量 $\Delta p_{x} = -2 m u_{1}$. By Newton's third law the wall gets an \textbf{impulse} 冲量 of $+2 m u_{1}$.

\item \textbf{time between hits on that wall}: travel $2L$ there and back, so $\Delta t = 2L/u_{1}$.

\item \textbf{average force from this molecule}: $F_{1} = \Delta p / \Delta t = m u_{1}^{2} / L$.

\item \textbf{add over all molecules}: $F = (Nm/L)\langle u_{x}^{2} \rangle$, where $\langle u_{x}^{2} \rangle$ is the \textbf{mean square} 均方 of the $x$-velocity.

\item \textbf{pressure}: $p = F/L^{2} = N m \langle u_{x}^{2} \rangle / V$.

\end{itemize}

In 3-D, by symmetry $\langle u_{x}^{2} \rangle = \tfrac{1}{3} \langle c^{2} \rangle$, where $\langle c^{2} \rangle$ is the \textbf{mean-square speed} 均方速率. So

\[
p V = \tfrac{1}{3} N m \langle c^{2} \rangle.
\]

\subsection*{Root-mean-square speed}

The square root of $\langle c^{2} \rangle$ is the \textbf{root-mean-square} 均方根 (r.m.s.) speed:

\[
c_{\text{r.m.s.}} = \sqrt{\langle c^{2} \rangle}.
\]

It is a useful single measure of how fast the molecules move, slightly larger than the mean speed (squaring weights fast molecules more).

\section*{Average translational kinetic energy}

Compare the two expressions for $pV$:

\[
p V = N k T \quad\text{and}\quad p V = \tfrac{1}{3} N m \langle c^{2} \rangle.
\]

Set them equal, cancel $N$, and multiply by $\tfrac{3}{2}$:

\[
\tfrac{3}{2} k T = \tfrac{1}{2} m \langle c^{2} \rangle.
\]

The right side is the \textbf{average translational kinetic energy} 平动动能 $\langle E_{\text{k}} \rangle$ of one molecule. So

\[
\langle E_{\text{k}} \rangle = \tfrac{1}{2} m \langle c^{2} \rangle = \tfrac{3}{2} k T.
\]

This is a key result: \textbf{the average translational kinetic energy of an ideal-gas molecule depends only on the thermodynamic temperature}, not on the type of gas or its pressure.

\subsection*{Consequences}

\begin{itemize}
\item \textbf{doubling the absolute temperature doubles the average KE} of each molecule, so $\langle c^{2} \rangle$ doubles and $c_{\text{r.m.s.}}$ grows by $\sqrt{2}$.

\item for \textbf{two gases at the same temperature}, the lighter gas has a larger $\langle c^{2} \rangle$. Hydrogen molecules move faster on average than oxygen molecules in the same room.

\item \textbf{total translational KE} of $N$ molecules: $\tfrac{3}{2} N k T = \tfrac{3}{2} n R T$.

\end{itemize}

\subsection*{Internal energy of an ideal gas}

For an ideal gas the molecules are point particles with no intermolecular potential energy and (in this simple model) no rotation or vibration. So the \textbf{internal energy} 内能 is just the total \textbf{kinetic energy} 动能:

\[
U = \tfrac{3}{2} N k T = \tfrac{3}{2} n R T.
\]

So the internal energy of an ideal gas is proportional to the thermodynamic temperature — doubling $T$ doubles $U$.

\subsection*{Changing pressure at fixed temperature}

Doubling $p$ at fixed $T$ (by squeezing the gas to half its volume) does \textbf{not} change $\langle E_{\text{k}} \rangle$ — that depends only on $T$. There are more wall collisions per second, but each molecule has the same average kinetic energy.


\end{document}
